% ======================================================================
% Topic 3 -- Counting and Probability -- part a (theory, examples, problems 1--15)
% ======================================================================
% helper: spinner with 6 equal sectors, labels clockwise starting from the top-left sector
\providecommand{\cpSpinSix}[6]{%
  \foreach \a in {90,150,210,270,330,30}{\draw[Ink,line width=0.8pt] (0,0)--(\a:1.1);}
  \draw[Ink,line width=1pt] (0,0) circle (1.1);
  \node[font=\itshape\large] at (120:0.72){#1};
  \node[font=\itshape\large] at (60:0.72){#2};
  \node[font=\itshape\large] at (0:0.72){#3};
  \node[font=\itshape\large] at (-60:0.72){#4};
  \node[font=\itshape\large] at (-120:0.72){#5};
  \node[font=\itshape\large] at (180:0.72){#6};
  \draw[-{Stealth[length=7pt]},line width=1.6pt,Ink] (0,0)--(250:0.98);
  \fill[Ink] (0,0) circle (1.6pt);}
% helper: a regular octagon centred at (#1,#2), width 1
\providecommand{\cpOct}[2]{%
  \draw[Ink,line width=0.8pt,fill=tmBlue!8] (#1-0.2071,#2+0.5)--(#1+0.2071,#2+0.5)--(#1+0.5,#2+0.2071)--(#1+0.5,#2-0.2071)
     --(#1+0.2071,#2-0.5)--(#1-0.2071,#2-0.5)--(#1-0.5,#2-0.2071)--(#1-0.5,#2+0.2071)--cycle;}
% helper: small square (diamond) between four octagons, centre (#1,#2)
\providecommand{\cpDia}[2]{%
  \draw[Ink,line width=0.8pt,fill=tmOrange!35] (#1,#2+0.2929)--(#1+0.2929,#2)--(#1,#2-0.2929)--(#1-0.2929,#2)--cycle;}
% helper: a flat marble (no 3D effect) centred at (#1,#2), colour #3
\providecommand{\rvcaBall}[3]{\filldraw[fill=#3!55,draw=#3!75!black,line width=0.7pt] (#1,#2) circle (0.22);}


\begin{knowbox}
\textbf{The big idea.} In this topic every problem comes back to \emph{one core skill}: \textbf{counting correctly} -- missing nothing and counting nothing twice.
\begin{itemize}[leftmargin=1.2em,itemsep=1pt,topsep=2pt]
\item \textbf{Counting} answers the question ``\emph{how many?}'': how many numbers, how many seating arrangements, how many rectangles in a figure, how many toothpicks\ldots
\item \textbf{Probability} answers the question ``\emph{how likely is it?}''. When all outcomes are equally likely, we only need to \emph{count twice}:
\[
\text{Probability}=\frac{\text{number of favourable outcomes}}{\text{number of all possible outcomes}} .
\]
\item Our counting tools: \textbf{systematic listing}, \textbf{tables}, \textbf{tree diagrams}, the \textbf{multiplication rule}, \textbf{splitting into cases} (the addition rule), \textbf{finding a pattern} and \textbf{reasoning by elimination}.
\end{itemize}
Always ask yourself three questions before starting: \emph{(1) How many possible outcomes are there altogether? (2) Are these outcomes ``fair'' (equally likely)? (3) How many of them satisfy the requirement?}
\end{knowbox}

\begin{figC}
\node[draw=tmBlue,fill=tmBlue!10,rounded corners=4pt,line width=1pt,minimum width=3.2cm,minimum height=0.9cm,font=\sffamily\bfseries] (D) at (0,0) {COUNT CORRECTLY};
\node[draw=tmGreen,fill=tmGreen!10,rounded corners=4pt,line width=1pt,minimum width=3.2cm,minimum height=0.9cm,font=\sffamily\bfseries] (X) at (6,0) {PROBABILITY};
\draw[->,line width=1.2pt,Ink] (D)--node[above,font=\footnotesize]{favourable / all}(X);
\foreach \t/\y in {Listing -- tables/1.9,Tree diagrams/1.2,Multiplication rule/0.5,Splitting into cases/-0.2,Finding a pattern/-0.9,Elimination/-1.6}{
  \node[draw=tmOrange,fill=tmOrange!10,rounded corners=3pt,font=\footnotesize,minimum width=2.9cm] at (-4.3,\y-0.15) (n) {\t};
  \draw[->,tmOrange!80!black] (n.east)--(D.west);}
\end{figC}
\vspace{0.3cm}

% ----------------------------------------------------------------------
\shead{Experiments, outcomes, events and probability}
\begin{knowbox}
\textbf{Key words.}
\begin{itemize}[leftmargin=1.2em,itemsep=1pt,topsep=2pt]
\item A \textbf{random experiment} is an action for which we know in advance \emph{all the possible outcomes}, but not \emph{which outcome will happen}. Examples: rolling a die, tossing a coin, spinning a spinner, drawing a marble from a bag without looking.
\item The \textbf{set of possible outcomes} is the list of \emph{all} outcomes. Rolling a die: $\{1,2,3,4,5,6\}$ -- $6$ outcomes. Tossing a coin: $\{\text{heads},\ \text{tails}\}$ -- $2$ outcomes.
\item An \textbf{event} is something we are interested in, for example ``rolling an even number''. The outcomes that make the event happen are called \textbf{favourable outcomes} for the event. The event ``rolling an even number'' has $3$ favourable outcomes: $2,4,6$.
\item Outcomes are \textbf{equally likely} if none of them is ``easier to get'' than another (a fair die, a fair coin, equal sectors on a spinner, marbles of identical shape and size).
\end{itemize}
\textbf{Probability formula} (when the outcomes are equally likely):
\[
P(\text{event})=\frac{\text{number of outcomes favourable to the event}}{\text{number of all possible outcomes}} .
\]
\textbf{Properties.} The number of favourable outcomes is between $0$ and the total number of outcomes, so always $0\le P\le 1$.
\begin{itemize}[leftmargin=1.2em,itemsep=1pt,topsep=2pt]
\item $P=0$: an \emph{impossible} event (for example rolling a $7$ on a die).
\item $P=1$: a \emph{certain} event (for example rolling a number less than $7$).
\end{itemize}
\end{knowbox}

\begin{proofbox}
Why a division? Suppose there are $n$ equally likely outcomes. The total ``chance'' $1$ (one of the outcomes is sure to happen) is shared \emph{equally} among $n$ outcomes, so each outcome has probability $\dfrac{1}{n}$. If the event has $k$ favourable outcomes, we add $k$ equal parts: $\dfrac1n+\dfrac1n+\dots+\dfrac1n=\dfrac{k}{n}$.
\end{proofbox}

\begin{figC}
\begin{scope}
\cpSpinSix{1}{2}{3}{4}{5}{6}
\node[font=\footnotesize,align=center] at (0,-1.55){Spinner with $6$ equal sectors:\\ each sector has probability $\frac16$};
\end{scope}
\begin{scope}[xshift=5cm]
\draw[Ink,line width=1pt,rounded corners=6pt] (-1.3,1)--(-1.2,-1)--(1.2,-1)--(1.3,1);
\foreach \x/\y/\c in {-0.75/-0.7/AnsRed,-0.25/-0.7/AnsRed,0.25/-0.7/StarGold,0.75/-0.7/StarGold,-0.5/-0.25/StarGold,
   0/-0.25/StarGold,0.5/-0.25/StarGold,-0.75/0.2/tmBlue,-0.25/0.2/tmBlue,0.25/0.2/tmBlue,0.75/0.2/tmBlue}{
  \rvcaBall{\x}{\y}{\c}}
\node[font=\footnotesize,align=center] at (0,-1.55){Bag with $2$ red, $5$ yellow, $4$ blue:\\ $P(\text{red})=\frac{2}{11}$};
\end{scope}
\end{figC}

\begin{warnbox}
\textbf{Count \emph{objects}, not \emph{kinds}.} The bag above has $3$ \emph{colours} but $11$ \emph{marbles}. The probability of drawing a red marble is $\dfrac{2}{11}$, \textbf{not} $\dfrac13$. Each marble is one equally likely outcome; the three colours are not equally likely (yellow has more marbles, so it is easier to draw).
\end{warnbox}

\begin{warnbox}
\textbf{Use the formula only when the outcomes are equally likely.} On a spinner with big and small sectors, the sectors are \emph{not} equally likely. Then the probability of landing in a sector equals \emph{the fraction of the whole circle} that the sector takes up, that is $\dfrac{\text{central angle of the sector}}{360^\circ}$. For example a sector with angle $90^\circ$ has probability $\dfrac{90}{360}=\dfrac14$.
\end{warnbox}

\begin{tipbox}
In multiple-choice questions the answer may be written in \emph{unsimplified} form, such as $\dfrac{2}{6}$ instead of $\dfrac13$. Compare the values of the fractions, not just how they look. On the other hand, when you write your own solution, simplify your final answer.
\end{tipbox}

% ----------------------------------------------------------------------
\shead{Complementary events. Finding a quantity from a probability}
\begin{knowbox}
\textbf{Complementary event.} The event ``$A$ does \emph{not} happen'' is called the complement of $A$. We have
\[
P(\text{not } A)=1-P(A).
\]
Example: rolling a die, $P(\text{a }6)=\dfrac16$, so $P(\text{not a }6)=1-\dfrac16=\dfrac56$.

\medskip
\textbf{Finding a quantity from a probability.} From the probability formula we get
\[
\text{number of favourable outcomes}=P\times\text{number of all outcomes}.
\]
Example: a bag has $20$ sweets and the probability of drawing a lemon sweet is $\dfrac{3}{4}$; then there are $20\times\dfrac34=15$ lemon sweets.

\medskip
\textbf{When the quantities change} (drawn and \emph{not replaced}, eaten, added\ldots): \textbf{recount both the numerator and the denominator} before computing the probability for the next draw.
\end{knowbox}

\begin{proofbox}
Suppose there are $n$ outcomes, of which $k$ are favourable to $A$. Each of the remaining $n-k$ outcomes makes $A$ \emph{not} happen. So
$P(\text{not }A)=\dfrac{n-k}{n}=\dfrac{n}{n}-\dfrac{k}{n}=1-P(A)$.
\end{proofbox}

\begin{warnbox}
A box has $10$ red and $10$ blue marbles. A student takes out $1$ blue marble and does \emph{not put it back}. What is the probability that the next marble is red?
\begin{itemize}[leftmargin=1.2em,itemsep=1pt,topsep=2pt]
\item \textbf{Wrong:} ``still $\dfrac{10}{20}$'' (forgetting that the box now has one marble fewer).
\item \textbf{Right:} $19$ marbles are left, and all $10$ red ones are still there, so the probability is $\dfrac{10}{19}$.
\end{itemize}
\end{warnbox}

\begin{parentbox}
Do a real experiment with sweets or bottle caps of different colours: put them in a cloth bag, count together ``how many are there altogether, how many are red'', then take one out and ask the same two questions again. Your child will see for themselves that both the numerator and the denominator change.
\end{parentbox}

% ----------------------------------------------------------------------
\shead{Systematic listing: lists, tables and tree diagrams}
\begin{knowbox}
When the number of possibilities is not too large, the safest method is to \textbf{list them all}. To list without missing or repeating anything, we need \textbf{an order}:
\begin{itemize}[leftmargin=1.2em,itemsep=1pt,topsep=2pt]
\item An \textbf{ordered list} (from smallest to largest, alphabetically\ldots). Example: the two-digit numbers with different digits made from $1,2,3$: $12,13,21,23,31,32$ -- ordered by the tens digit.
\item A \textbf{two-way table} when the experiment has \emph{two stages} (rolling two dice, choosing a main dish and a dessert\ldots): each row is an outcome of stage 1, each column an outcome of stage 2, each cell a combined outcome.
\item A \textbf{tree diagram} when there are \emph{two, three, four stages in a row}: each stage branches out; each ``path'' from the root to the end of a branch is one outcome.
\end{itemize}
\end{knowbox}

\needspace{14\baselineskip}
\textbf{Table for rolling two dice (one red, one blue).} Each cell shows the \emph{sum} of the two numbers. There are $6$ rows and $6$ columns, so $6\times 6=36$ equally likely outcomes.
\begin{figC}
\node[font=\bfseries\footnotesize,text=AnsRed] at (-0.15,0.75) {red $\downarrow$};
\node[font=\bfseries\footnotesize,text=tmBlue] at (2.1,1.05) {blue $\rightarrow$};
\foreach \c in {1,...,6}{\node[draw=tmBlue,fill=tmBlue!15,minimum size=0.62cm,font=\bfseries] at (\c*0.62,0.4) {\c};}
\foreach \r in {1,...,6}{\node[draw=AnsRed,fill=AnsRed!12,minimum size=0.62cm,font=\bfseries] at (0,0.4-\r*0.62) {\r};
  \foreach \c in {1,...,6}{\pgfmathtruncatemacro\s{\r+\c}%
    \ifnum\s=7 \def\fc{tmOrange!35}\else\def\fc{white}\fi
    \node[draw=Line,fill=\fc,minimum size=0.62cm] at (\c*0.62,0.4-\r*0.62) {\s};}}
\node[font=\footnotesize,align=left,anchor=west] at (4.6,-1.7) {Orange cells: sum equals $7$.\\ There are $6$ such cells, so\\ $P(\text{sum}=7)=\dfrac{6}{36}=\dfrac16$.};
\end{figC}

\begin{warnbox}
\textbf{(1, 2) and (2, 1) are two different outcomes!} Red $1$, blue $2$ is different from red $2$, blue $1$. If you only count ``pairs of numbers'' without order, the pairs are \emph{no longer} equally likely (the pair $\{1,2\}$ can happen in $2$ ways, the pair $\{1,1\}$ in only $1$ way). That is why the sums $2,3,\dots,12$ are \emph{not} equally likely: a sum of $7$ can happen in $6$ ways, a sum of $2$ in only $1$ way.
\end{warnbox}

\textbf{Tree diagram for tossing a coin three times} (H: heads, T: tails). Each toss has $2$ branches, so there are $2\times2\times2=8$ paths -- that is, $8$ equally likely outcomes.
\begin{center}\begin{tikzpicture}[font=\sffamily\small,grow'=right,level distance=2.1cm,
  every node/.style={circle,draw=tmBlue,fill=tmBlue!10,inner sep=1.5pt,minimum size=0.55cm,font=\footnotesize\bfseries},
  level 1/.style={sibling distance=2.6cm},level 2/.style={sibling distance=1.3cm},level 3/.style={sibling distance=0.65cm},
  edge from parent/.style={draw=Ink!70,line width=0.7pt}]
\node[fill=Ink!10,draw=Ink] {Start}
  child {node {H}
    child {node {H} child {node {H}} child {node {T}}}
    child {node {T} child {node {H}} child {node {T}}}}
  child {node {T}
    child {node {H} child {node {H}} child {node {T}}}
    child {node {T} child {node {H}} child {node {T}}}};
\foreach \y/\t in {2.275/HHH,1.625/HHT,0.975/HTH,0.325/HTT,-0.325/THH,-0.975/THT,-1.625/TTH,-2.275/TTT}{
  \node[draw=none,fill=none,font=\footnotesize,anchor=west] at (6.7,\y) {\t};}
\node[draw=none,fill=none,font=\footnotesize] at (2.1,-2.8) {toss 1};
\node[draw=none,fill=none,font=\footnotesize] at (4.2,-2.8) {toss 2};
\node[draw=none,fill=none,font=\footnotesize] at (6.3,-2.8) {toss 3};
\end{tikzpicture}\end{center}
Example: the event ``exactly two heads'' consists of HHT, HTH, THH -- $3$ outcomes, so its probability is $\dfrac38$.

\begin{tipbox}
A problem like ``walk on a grid by tossing a coin'' (heads: move up, tails: move right) is really a problem about counting sequences of heads and tails: after $4$ tosses there are $2^4=16$ equally likely sequences, and the \emph{end point} depends only on the \emph{number of heads}, not on their order.
\end{tipbox}

% ----------------------------------------------------------------------
\shead{The multiplication rule and the addition rule}
\begin{knowbox}
\textbf{Multiplication rule.} If a task consists of two successive stages, stage 1 can be done in $m$ ways, and \emph{for each} way of doing stage 1, stage 2 can be done in $n$ ways, then the whole task can be done in $m\times n$ ways. (With three or four stages: keep multiplying.)

\medskip
\textbf{Addition rule (splitting into cases).} If the ways are split into cases that \emph{do not overlap}: case 1 has $a$ ways, case 2 has $b$ ways, \ldots, then the total number of ways is $a+b+\dots$

\medskip
\textbf{Common situations with digits.}
\begin{itemize}[leftmargin=1.2em,itemsep=1pt,topsep=2pt]
\item Three-digit numbers made from $3$ given non-zero digits, \emph{no repetition}: $3\times2\times1=6$ numbers.
\item Three-digit numbers made from $3$ given non-zero digits, \emph{repetition allowed}: $3\times3\times3=27$ numbers.
\item The first (leading) digit \textbf{cannot be $0$}: $9$ choices ($1$ to $9$); the other places have $10$ choices ($0$ to $9$).
\item A three-digit \textbf{palindrome} (reads the same forwards and backwards) has the form $\overline{aba}$: $9$ ways to choose $a$, $10$ ways to choose $b$, and the last digit is ``forced'' to equal $a$ (only $1$ way). There are $9\times10=90$ such numbers.
\end{itemize}
\end{knowbox}

\countboxes{hundreds/9/from 1 to 9/tmBlue,tens/10/from 0 to 9/tmOrange,ones/1/same as hundreds/tmGreen}{90}
\vspace{-6pt}
\begin{center}\footnotesize Number of three-digit palindromes: $9\times 10\times 1=90$.\end{center}

\begin{proofbox}
Why is the multiplication rule true? Draw a table with $m$ rows (the ways of doing stage 1) and $n$ columns (the ways of doing stage 2). Each cell of the table is \emph{exactly one} way of doing the whole task, and each way of doing the whole task is in \emph{exactly one} cell. The table has $m$ rows of $n$ cells, so $m\times n$ cells. A tree diagram gives the same result: $m$ big branches, each splitting into $n$ small branches.
\end{proofbox}

\begin{knowbox}
\textbf{Choosing two objects (order does not matter).} Choosing $2$ of $n$ students for classroom duty: if we care about order (first person, second person) there are $n\times(n-1)$ ways; but each pair is counted \emph{twice} (Anna--Ben and Ben--Anna), so the number of pairs is
\[
\frac{n\times(n-1)}{2}.
\]
Example: there are $\dfrac{10\times9}{2}=45$ pairs of different numbers taken from $10$ numbers.
\end{knowbox}

\begin{warnbox}
\textbf{Double counting.} When the question asks ``how many \emph{pairs}'', ``how many \emph{groups}'' (who comes first does not matter) and you use the multiplication rule with order, each group is counted several times -- you must divide. On the other hand, when the question is about \emph{numbers} (the order of the digits matters) or \emph{seating in a row}, order matters and you do \emph{not} divide.
\end{warnbox}

\begin{tipbox}
When splitting into cases, ask: ``Do the cases \emph{cover} every possibility? Do they \emph{overlap}?'' For example, for the sum $m+n$ to be even, $m$ and $n$ must be \emph{both odd} or \emph{both even} -- these two cases are separate and together they cover every possibility.
\end{tipbox}

% ----------------------------------------------------------------------
\shead{Reasoning by elimination: seating and colouring}
\begin{knowbox}
Some problems do not ask ``how many'' but ``\emph{who can}, \emph{who cannot}'', ``\emph{what is the least number needed}''. Then we use \textbf{reasoning}:
\begin{itemize}[leftmargin=1.2em,itemsep=1pt,topsep=2pt]
\item \textbf{Draw a diagram} of the positions (seats, cells, faces\ldots) and try filling it in.
\item Two people who \textbf{must sit together} can be treated as \textbf{one block} taking two neighbouring seats.
\item To show ``\emph{X can} sit in the middle seat'': just \textbf{give one arrangement} that satisfies all the conditions.
\item To show ``\emph{X cannot} sit in the middle seat'': \textbf{suppose} X sits in the middle seat and show that \emph{every} remaining arrangement breaks a condition (leading to a contradiction). This is called \textbf{proof by contradiction}.
\item A ``what is the \emph{least} number'' problem has \textbf{two parts}: (1) show a way that works with that number; (2) explain why fewer is impossible.
\end{itemize}
\end{knowbox}

\begin{figC}
\foreach \i in {1,...,5}{\draw[tmBlue,line width=1pt,fill=tmBlue!8,rounded corners=3pt] (\i*1.3-0.5,0) rectangle (\i*1.3+0.5,0.7);
  \node[font=\footnotesize] at (\i*1.3,-0.3){seat \i};}
\draw[tmOrange,line width=1.4pt,rounded corners=5pt] (3*1.3-0.62,-0.12) rectangle (3*1.3+0.62,0.82);
\node[font=\footnotesize\bfseries,text=tmOrange!80!black] at (3*1.3,1.1){middle seat};
\draw[decorate,decoration={brace,amplitude=5pt,mirror},Ink] (1*1.3-0.5,-0.55)--(2*1.3+0.5,-0.55) node[midway,below=5pt,font=\footnotesize]{two seats on the left};
\draw[decorate,decoration={brace,amplitude=5pt,mirror},Ink] (4*1.3-0.5,-0.55)--(5*1.3+0.5,-0.55) node[midway,below=5pt,font=\footnotesize]{two seats on the right};
\end{figC}
Key observation: once someone sits in the middle seat, the other four seats split into \emph{two pairs of neighbouring seats} (seats 1--2 and seats 4--5). The two seats on either side of the middle seat (seat 2 and seat 4) are \emph{not} next to each other.

\begin{knowbox}
\textbf{Colouring so that neighbouring parts have different colours.} Two faces of a cube are ``neighbours'' if they share an edge. Each face is a neighbour of $4$ faces and is \emph{not} a neighbour of exactly $1$ face -- the \textbf{opposite} face. A cube has $3$ pairs of opposite faces (top--bottom, front--back, left--right).
\begin{itemize}[leftmargin=1.2em,itemsep=1pt,topsep=2pt]
\item \emph{$3$ colours are enough:} paint each pair of opposite faces one colour; two faces of the same colour are opposite, so they are not neighbours.
\item \emph{Fewer than $3$ colours is impossible:} the three faces meeting at a corner are neighbours of each other in pairs, so they need three different colours.
\end{itemize}
\end{knowbox}

\begin{tipbox}
For a multiple-choice question ``Who \emph{cannot}\ldots?'', you can work backwards: for each option, try to find \emph{one} valid arrangement. Every option for which you find an arrangement is eliminated; the remaining option is the answer -- but remember to explain why it really cannot be arranged.
\end{tipbox}

% ----------------------------------------------------------------------
\shead{Counting shapes in a figure}
\begin{knowbox}
\textbf{Principle.} Counting shapes is a listing problem. To miss nothing and count nothing twice, \textbf{split into types} (by size, by position) and then \textbf{add up}.

\medskip
\textbf{The ``choose two lines'' trick.} A rectangle in a grid is completely determined by \emph{two horizontal lines} (top and bottom sides) and \emph{two vertical lines} (left and right sides). Therefore:
\begin{itemize}[leftmargin=1.2em,itemsep=1pt,topsep=2pt]
\item A strip of $1$ row and $n$ squares has $n+1$ vertical lines; the number of rectangles is the number of ways to choose $2$ vertical lines:
$\dfrac{(n+1)\times n}{2}$. Example: a strip of $4$ squares has $\dfrac{5\times4}{2}=10$ rectangles.
\item A grid with $a$ horizontal lines and $b$ vertical lines (all running the full length) has
$\dfrac{a(a-1)}{2}\times\dfrac{b(b-1)}{2}$ rectangles (multiplication rule!).
\item The same counting works for \emph{parallelograms} in a grid whose slanted lines are parallel to each other.
\end{itemize}
\end{knowbox}

\begin{figC}
\begin{scope}
\fill[tmOrange!30] (1,0) rectangle (3,1);
\foreach \i in {0,...,4}{\draw[Ink,line width=1pt] (\i,0)--(\i,1); \node[font=\footnotesize,text=tmBlue] at (\i,-0.3){\pgfmathtruncatemacro\k{\i+1}\k};}
\draw[Ink,line width=1pt] (0,0)--(4,0) (0,1)--(4,1);
\node[font=\footnotesize,align=center] at (2,-0.95){Choose vertical lines $2$ and $4$\\ $\Rightarrow$ the orange rectangle};
\end{scope}
\begin{scope}[xshift=6cm]
\fill[tmGreen!30] (1,0) rectangle (3,2);
\foreach \x in {0,1,2,3}{\draw[Ink,line width=1pt] (\x,0)--(\x,2);}
\foreach \y in {0,1,2}{\draw[Ink,line width=1pt] (0,\y)--(3,\y);}
\node[font=\footnotesize,align=center] at (1.5,-0.7){$3$ horizontal, $4$ vertical lines:\\ $3\times 6=18$ rectangles};
\end{scope}
\end{figC}

\begin{proofbox}
Why does a strip of $n$ squares have $\dfrac{(n+1)\times n}{2}$ rectangles? Choose the left side: $n+1$ ways; choose the right side (different from the left side): $n$ ways -- that gives $(n+1)\times n$ \emph{ordered} choices. But choosing ``line 2 then line 4'' or ``line 4 then line 2'' gives the same rectangle, so each rectangle is counted $2$ times; divide by $2$. Counting by size gives the same result: a strip of $4$ squares has $4$ rectangles of width $1$, $3$ of width $2$, $2$ of width $3$ and $1$ of width $4$, in total $4+3+2+1=10$.
\end{proofbox}

\begin{warnbox}
When the shapes \textbf{partly overlap} (not a grid with full-length lines), the ``choose two lines'' trick \emph{does not work}, because many segments do not run all the way across. Check each shape: \textbf{are all four sides really drawn?} Do not count ``imaginary rectangles'' that are missing a side.
\end{warnbox}

\begin{knowbox}
\textbf{Counting in a tiling.} A tiling made of regular octagons and small squares: if there are $m$ rows with $n$ octagons in each row, there are $m\times n$ octagons. Each small square sits where \emph{four octagons meet}; between $n$ columns there are $n-1$ gaps and between $m$ rows there are $m-1$ gaps, so there are $(m-1)\times(n-1)$ squares.
\end{knowbox}
\begin{figC}
\foreach \i in {0,1,2}{\foreach \j in {0,1}{\cpOct{\i}{\j}}}
\foreach \i in {0.5,1.5}{\cpDia{\i}{0.5}}
\node[font=\footnotesize,align=left,anchor=west] at (3,0.5){$2$ rows, $3$ columns: $2\times3=6$ octagons,\\ $(2-1)\times(3-1)=2$ squares.};
\end{figC}

% ----------------------------------------------------------------------
\shead{Counting with patterns: evenly spaced sequences, constant rates}
\begin{knowbox}
\textbf{Counting the whole numbers from $a$ to $b$} (including $a$ and $b$): there are $b-a+1$ numbers. For example from $1$ to $11$ there are $11$ numbers, of which $6$ are odd ($1,3,5,7,9,11$) and $5$ are even.

\medskip
\textbf{Evenly spaced sequence} (each number is a fixed amount $d$ more than the one before):
\[
\text{number of terms}=\frac{\text{last}-\text{first}}{d}+1,\qquad
\text{sum}=\frac{(\text{first}+\text{last})\times\text{number of terms}}{2}.
\]

\textbf{Quantities that change at a constant rate} (moving steadily, flowing steadily, counting steadily\ldots):
$\text{amount}=\text{rate}\times\text{time}$. For example, if $3$ sheep pass each minute, then in $14$ minutes $3\times14=42$ sheep pass.

\medskip
\textbf{Figures made of sticks.} Count the \emph{horizontal sticks} and the \emph{vertical sticks} separately and then add -- usually easier than counting square by square.

\medskip
\textbf{Events that repeat in cycles.} If one thing repeats every $a$ minutes and another every $b$ minutes, and they start together, they coincide again every $\text{LCM}(a,b)$ minutes.
\end{knowbox}

\begin{proofbox}
\textbf{Gauss's way of adding.} Compute $S=1+2+3+\dots+10$. Write the sum twice, once forwards and once backwards:
\[
\begin{array}{rcccccc}
S&=&1&+\;2&+\;3&+\dots+&10\\
S&=&10&+\;9&+\;8&+\dots+&1\\ \hline
2S&=&11&+\;11&+\;11&+\dots+&11
\end{array}
\]
There are $10$ columns, each with sum $11$, so $2S=10\times11=110$ and $S=55$. In general: sum $=\dfrac{(\text{first}+\text{last})\times\text{number of terms}}{2}$.
\end{proofbox}

\begin{warnbox}
\textbf{The ``off by one'' error.} From $3$ to $10$ there are $10-3+1=8$ numbers, \emph{not} $7$. A fence $10$ m long with a post every $1$ m (including both ends) needs $11$ posts, not $10$. Test a formula on a small example to check it.
\end{warnbox}

\begin{parentbox}
For ``toothpick'' problems, let your child actually build the small figures ($1$ square, $2$ squares, $3$ squares\ldots) with toothpicks, write the number of toothpicks in a table, and then ask: ``each time we add one square, how many toothpicks do we add?''. A child who finds the pattern alone remembers it much longer than a memorised formula.
\end{parentbox}

\needspace{16\baselineskip}
\begin{knowbox}
\textbf{Summary: read the clues, choose the tool.}\par\smallskip
\begin{tabularx}{\linewidth}{@{}>{\raggedright\arraybackslash}X>{\raggedright\arraybackslash}X@{}}
\textbf{Clue in the problem} & \textbf{Tool to use}\\ \hline
``choose one at random'', ``what is the probability'' & count all outcomes and favourable outcomes, then divide\\
probability given, quantity asked & favourable $=P\times$ total\\
something was taken out/eaten before the next draw & recount numerator and denominator\\
``at least one'', ``not'' & complement: $1-P$\\
two dice, two independent choices & two-way table\\
tossing a coin several times, several stages in a row & tree diagram\\
``how many numbers can be made'', ``how many ways'' & multiplication rule (choose restricted positions first)\\
``how many pairs'' (order does not matter) & multiplication rule, then divide by $2$\\
conditions of the form ``either\ldots or\ldots'' & split into cases, add\\
``who can / cannot'' & try one arrangement / proof by contradiction\\
``how many shapes\ldots'' & split by size, or choose two lines\\
rows of seats, toothpicks, constant rates & find the pattern, evenly spaced sequences\\
\end{tabularx}
\end{knowbox}

% ======================================================================
% EXAMPLES
% ======================================================================
\begin{exbox}{Example 1}
A box contains $3$ red marbles, $5$ blue marbles and $4$ yellow marbles, all of the same size. One marble is drawn at random.
\begin{enumerate}[label=\alph*),leftmargin=1.5em,itemsep=0pt,topsep=2pt]
\item Find the probability of drawing a blue marble.
\item Find the probability of drawing a marble that is \emph{not} red.
\item Find the probability of drawing a purple marble.
\end{enumerate}
\guide
\step{1}{Count all possible outcomes}
Each marble is one outcome, and since the marbles have the same size the outcomes are equally likely. Total number of marbles: $3+5+4=12$. So there are $12$ possible outcomes.
\step{2}{Part a}
There are $5$ blue marbles, i.e.\ $5$ favourable outcomes.
$P(\text{blue})=\dfrac{5}{12}$.
\step{3}{Part b -- use the complement}
$P(\text{red})=\dfrac{3}{12}=\dfrac14$, so $P(\text{not red})=1-\dfrac14=\dfrac34$.

\emph{Another way (direct counting):} the marbles that are not red are $5$ blue and $4$ yellow, i.e.\ $9$ marbles; $P=\dfrac{9}{12}=\dfrac34$. Both ways give the same result.
\step{4}{Part c}
There are no purple marbles in the box: $0$ favourable outcomes. $P(\text{purple})=\dfrac{0}{12}=0$ -- an impossible event.
\ans{a) $\frac{5}{12}$;\quad b) $\frac34$;\quad c) $0$.}
\end{exbox}

\begin{exbox}{Example 2}
A bag contains $30$ sweets, some strawberry and some orange. If one sweet is drawn at random, the probability that it is strawberry is $\dfrac25$.
\begin{enumerate}[label=\alph*),leftmargin=1.5em,itemsep=0pt,topsep=2pt]
\item How many strawberry sweets and how many orange sweets are in the bag?
\item Anna takes out $2$ strawberry sweets and eats them. Then Ben draws one sweet at random from the bag. Find the probability that Ben gets a strawberry sweet.
\end{enumerate}
\guide
\step{1}{Find the number of strawberry sweets}
\begin{recallbox}Number of favourable outcomes $=$ probability $\times$ number of all outcomes.\end{recallbox}
Strawberry sweets: $30\times\dfrac25=\dfrac{30\times 2}{5}=\dfrac{60}{5}=12$. Orange sweets: $30-12=18$.

\emph{Check:} $\dfrac{12}{30}=\dfrac{2}{5}$ \checkmark.
\step{2}{Recount after Anna eats $2$ sweets}
\begin{center}
\begin{tabular}{|l|c|c|c|}\hline
 & strawberry & orange & total \\ \hline
At the start & $12$ & $18$ & $30$\\ \hline
After Anna eats $2$ strawberry & $12-2=10$ & $18$ & $30-2=28$\\ \hline
\end{tabular}
\end{center}
\step{3}{Compute the probability for Ben}
$P(\text{strawberry})=\dfrac{10}{28}=\dfrac{5}{14}$.
\begin{warnbox}Do not answer $\dfrac{10}{30}$: Anna has eaten $2$ sweets, so the total number of sweets (the denominator) has gone down too.\end{warnbox}
\ans{a) $12$ strawberry, $18$ orange;\quad b) $\frac{5}{14}$.}
\end{exbox}

\begin{exbox}{Example 3}
A red die and a blue die (both fair) are rolled together. Find the probability of each event:
\begin{enumerate}[label=\alph*),leftmargin=1.5em,itemsep=0pt,topsep=2pt]
\item The sum of the two numbers is $8$.
\item Both dice show the same number.
\item The sum of the two numbers is greater than or equal to $10$.
\end{enumerate}
\guide
\step{1}{Make a table of all outcomes}
Rows are the numbers on the red die, columns the numbers on the blue die, and each cell shows the sum. The table has $6\times6=36$ cells, corresponding to $36$ equally likely outcomes.
\begin{figC}
\foreach \c in {1,...,6}{\node[draw=tmBlue,fill=tmBlue!15,minimum size=0.6cm,font=\bfseries] at (\c*0.6,0.4) {\c};}
\foreach \r in {1,...,6}{\node[draw=AnsRed,fill=AnsRed!12,minimum size=0.6cm,font=\bfseries] at (0,0.4-\r*0.6) {\r};
  \foreach \c in {1,...,6}{\pgfmathtruncatemacro\s{\r+\c}%
    \def\fc{white}%
    \ifnum\s=8 \def\fc{tmOrange!40}\fi
    \ifnum\s>9 \def\fc{tmGreen!30}\fi
    \node[draw=Line,fill=\fc,minimum size=0.6cm] at (\c*0.6,0.4-\r*0.6) {\s};
    \ifnum\r=\c \draw[tmPurple,line width=1pt] (\c*0.6,0.4-\r*0.6) circle (0.24);\fi}}
\node[font=\footnotesize,align=left,anchor=west] at (4.4,-1.7) {\colorbox{tmOrange!40}{\strut\ \ } sum equals $8$ ($5$ cells)\\[2pt]
  {\color{tmPurple}$\bigcirc$} both numbers equal ($6$ cells)\\[2pt] \colorbox{tmGreen!30}{\strut\ \ } sum $\ge 10$ ($6$ cells)};
\end{figC}
\step{2}{Part a}
The cells with sum $8$: $(2,6),(3,5),(4,4),(5,3),(6,2)$ -- $5$ outcomes (the first number is the red die). $P=\dfrac{5}{36}$.
\step{3}{Part b}
Equal numbers lie on the diagonal: $(1,1),(2,2),\dots,(6,6)$ -- $6$ outcomes. $P=\dfrac{6}{36}=\dfrac16$.
\step{4}{Part c}
Sum $10$: $(4,6),(5,5),(6,4)$; sum $11$: $(5,6),(6,5)$; sum $12$: $(6,6)$. Altogether $3+2+1=6$ outcomes. $P=\dfrac{6}{36}=\dfrac16$.
\begin{tipbox}The $6\times6$ table also shows \emph{symmetry}: the number of cells above the diagonal (red smaller than blue) equals the number of cells below the diagonal (red greater than blue).\end{tipbox}
\ans{a) $\frac{5}{36}$;\quad b) $\frac16$;\quad c) $\frac16$.}
\end{exbox}

\needspace{26\baselineskip}
\begin{exbox}{Example 4}
A fair coin is tossed three times in a row. Find the probability of getting:
a) exactly two heads;\quad b) at least one head.
\guide
\step{1}{Draw a tree diagram and list}
Each toss has $2$ outcomes (H or T). By the tree diagram in section 3, there are $2\times2\times2=8$ equally likely outcomes:
\[
\text{HHH},\ \text{HHT},\ \text{HTH},\ \text{HTT},\ \text{THH},\ \text{THT},\ \text{TTH},\ \text{TTT}.
\]
\step{2}{Part a}
The outcomes with exactly two H's: HHT, HTH, THH -- $3$ outcomes. $P=\dfrac38$.
\step{3}{Part b -- complement}
The complement of ``at least one head'' is ``no heads at all'', which is only TTT ($1$ outcome). So
$P(\text{at least one H})=1-\dfrac18=\dfrac78$.
\begin{tipbox}When you see the words ``\emph{at least one}'', think immediately of the complement ``\emph{none at all}'' -- it usually has very few outcomes and is easy to count.\end{tipbox}
\ans{a) $\frac38$;\quad b) $\frac78$.}
\end{exbox}

\begin{exbox}{Example 5}
Using the digits $1,2,3,4$, how many three-digit numbers can be made if:
a) digits may be repeated;\quad b) the digits must be different;\quad c) the digits are different and the number is even?\\
d) One of the numbers from part b) is chosen at random. Find the probability that it is even.
\guide
\step{1}{Part a -- repetition allowed}
Hundreds: $4$ ways; tens: still $4$ ways; ones: $4$ ways. By the multiplication rule: $4\times4\times4=64$ numbers.
\step{2}{Part b -- no repetition}
\countboxes{hundreds/4/any/tmBlue,tens/3/not the hundreds digit/tmOrange,ones/2/not the two above/tmGreen}{24}
Hundreds: $4$ ways. The tens digit must differ from the digit already used: $3$ ways. The ones digit must differ from both digits used: $2$ ways. There are $4\times3\times2=24$ numbers.
\step{3}{Part c -- even, no repetition: choose the ones digit FIRST}
An even number ends in $2$ or $4$: $2$ ways. Then the hundreds digit is one of the $3$ remaining digits: $3$ ways; the tens digit: $2$ ways. There are $2\times3\times2=12$ numbers.
\begin{tipbox}A position with a \emph{condition} (must be even, must not be $0$\ldots) should be \textbf{chosen first}. If you choose the hundreds digit first, the number of choices for the ones digit is sometimes $1$ and sometimes $2$ (depending on whether the hundreds digit used an even or an odd digit) -- very easy to get wrong.\end{tipbox}
\step{4}{Part d}
$P(\text{even})=\dfrac{12}{24}=\dfrac12$.
\verify{List the numbers starting with $1$: $123,124,132,134,142,143$ -- exactly $6=3\times2$ numbers, of which $124,132,134,142$ are even ($4$ numbers). The four possible first digits behave alike, so there are $4\times6=24$ numbers, \ldots\ matching part b.}
\ans{a) $64$;\quad b) $24$;\quad c) $12$;\quad d) $\frac12$.}
\end{exbox}

\begin{exbox}{Example 6}
Four friends Anna, Ben, Chloe and David sit on four chairs in a row, numbered $1,2,3,4$ from left to right. We know that: David sits on chair $1$; Ben and Chloe sit next to each other; Anna does not sit next to Ben; Anna does not sit on chair $4$. Who sits on chair $4$?
\guide
\step{1}{Place David}
David is on chair $1$, leaving chairs $2,3,4$ for Anna, Ben and Chloe.
\step{2}{Consider Anna's position}
Anna does not sit on chair $4$, so Anna sits on chair $2$ or chair $3$.
\begin{itemize}[leftmargin=1.2em,itemsep=1pt,topsep=1pt]
\item \emph{If Anna sits on chair $3$:} Ben and Chloe take chairs $2$ and $4$ -- these chairs are not next to each other, which breaks the condition ``Ben and Chloe sit next to each other''. Rejected.
\item \emph{If Anna sits on chair $2$:} Ben and Chloe take chairs $3$ and $4$ (next to each other, good). Anna must not be next to Ben, and chair $3$ is next to chair $2$, so Ben is not on chair $3$. So Ben is on chair $4$ and Chloe on chair $3$.
\end{itemize}
\begin{figC}
\foreach \i/\n/\c in {1/David/tmPurple,2/Anna/tmOrange,3/Chloe/tmGreen,4/Ben/tmBlue}{
  \draw[\c,line width=1pt,fill=\c!12,rounded corners=3pt] (\i*1.6-0.65,0) rectangle (\i*1.6+0.65,0.7);
  \node[font=\bfseries] at (\i*1.6,0.35){\n};\node[font=\footnotesize] at (\i*1.6,-0.3){chair \i};}
\end{figC}
\verify{David on chair $1$ \checkmark; Ben (chair $4$) next to Chloe (chair $3$) \checkmark; Anna (chair $2$) is next to David and Chloe, not next to Ben \checkmark; Anna is not on chair $4$ \checkmark.}
\ans{Ben sits on chair $4$.}
\end{exbox}

\needspace{14\baselineskip}
\begin{exbox}{Example 7}
a) The figure on the left is a strip of $4$ squares. How many rectangles (including squares) are there altogether?\\
b) The figure on the right is a $2\times2$ grid of squares. How many rectangles are there? How many of them are squares?
\begin{figC}
\foreach \i in {0,...,4}{\draw[Ink,line width=1pt] (\i*0.8,0)--(\i*0.8,0.8);}
\draw[Ink,line width=1pt] (0,0)--(3.2,0) (0,0.8)--(3.2,0.8);
\begin{scope}[xshift=5.5cm]
\foreach \i in {0,1,2}{\draw[Ink,line width=1pt] (\i*0.8,0)--(\i*0.8,1.6) (0,\i*0.8)--(1.6,\i*0.8);}
\end{scope}
\end{figC}
\guide
\step{1}{Part a -- count by width}
Measure the width in squares:
\begin{center}
\begin{tabular}{|c|c|c|c|c|}\hline
Width & $1$ square & $2$ squares & $3$ squares & $4$ squares\\ \hline
Number of shapes & $4$ & $3$ & $2$ & $1$\\ \hline
\end{tabular}
\end{center}
Total: $4+3+2+1=10$ rectangles. \emph{Check with the ``choose two vertical lines'' trick:} there are $5$ vertical lines, and the number of ways to choose $2$ of them is $\dfrac{5\times4}{2}=10$ \checkmark.
\step{2}{Part b -- choose two horizontal and two vertical lines}
The grid has $3$ horizontal and $3$ vertical lines. Choose $2$ horizontal lines: $\dfrac{3\times2}{2}=3$ ways; choose $2$ vertical lines: $3$ ways. By the multiplication rule there are $3\times3=9$ rectangles.
\step{3}{Count the squares separately}
$1\times1$ squares: $4$; $2\times2$ squares: $1$ (the whole grid). There are $4+1=5$ squares.
\verify{Counting by size: $1\times1$: $4$; $1\times2$ (lying flat): $2$; $2\times1$ (standing): $2$; $2\times2$: $1$. Total $4+2+2+1=9$ \checkmark.}
\ans{a) $10$ rectangles;\quad b) $9$ rectangles, of which $5$ are squares.}
\end{exbox}

\begin{exbox}{Example 8}
Toothpicks are used to make a row of joined squares as shown below.
\begin{figC}
\foreach \n/\x in {1/0,2/2,3/5}{
  \foreach \i in {0,...,\n}{\draw[tmOrange!80!black,line width=1.6pt,line cap=round] (\x+\i*0.7,0.05)--(\x+\i*0.7,0.65);}
  \foreach \i in {1,...,\n}{\draw[tmOrange!80!black,line width=1.6pt,line cap=round] (\x+\i*0.7-0.65,0)--(\x+\i*0.7-0.05,0) (\x+\i*0.7-0.65,0.7)--(\x+\i*0.7-0.05,0.7);}
  \node[font=\footnotesize] at (\x+\n*0.35,-0.35){\n\ sq.};}
\end{figC}
a) How many toothpicks are needed for a row of $20$ squares?\\
b) Lucy builds, one after another and separately, a row of $1$ square, a row of $2$ squares, \ldots, a row of $10$ squares. How many toothpicks does Lucy use altogether?
\guide
\step{1}{Find the pattern}
\begin{center}
\begin{tabular}{|c|c|c|c|c|}\hline
Squares & $1$ & $2$ & $3$ & $n$\\ \hline
Toothpicks & $4$ & $7$ & $10$ & $3n+1$\\ \hline
\end{tabular}
\end{center}
Explanation: count each kind of toothpick separately. Horizontal toothpicks: each square has $1$ on top and $1$ at the bottom, that is $2n$ toothpicks. Vertical toothpicks: $n$ squares have $n+1$ vertical toothpicks (just as $n$ gaps need $n+1$ posts). Total $2n+(n+1)=3n+1$.
\step{2}{Part a}
$n=20$: $3\times20+1=61$ toothpicks.
\step{3}{Part b -- sum of an evenly spaced sequence}
The numbers of toothpicks are $4,7,10,\dots,31$ (a row of $10$ squares needs $3\times10+1=31$ toothpicks). This is an evenly spaced sequence with step $3$ and $10$ terms.
\begin{recallbox}Sum of an evenly spaced sequence $=\dfrac{(\text{first}+\text{last})\times\text{number of terms}}{2}$.\end{recallbox}
Sum: $\dfrac{(4+31)\times10}{2}=\dfrac{350}{2}=175$ toothpicks.
\verify{$4+7+10+13+16+19+22+25+28+31$: pair first with last, $4+31=7+28=10+25=13+22=16+19=35$, giving $5$ pairs, $5\times35=175$ \checkmark.}
\ans{a) $61$ toothpicks;\quad b) $175$ toothpicks.}
\end{exbox}

% ======================================================================
\clearpage
\dividerbar{Core practice}

% ---------------------------------------------------------------- Problem 1
\begin{probbox}{1}{\relax}
\begin{minipage}[t]{0.66\linewidth}\vspace{0pt}
The spinner shown is divided into $6$ sectors of equal size, each marked with a letter. The spinner is spun once. What is the probability that the arrow lands on a sector marked with the letter $P$?
\choices{$\dfrac{3}{6}$}{$\dfrac{4}{6}$}{$\dfrac{5}{6}$}{$\dfrac{2}{6}$}{$\dfrac{1}{6}$}
\end{minipage}\hfill
\begin{minipage}[t]{0.3\linewidth}\vspace{0pt}
\centering\begin{tikzpicture}[scale=0.95]\cpSpinSix{P}{R}{P}{Q}{R}{Q}\end{tikzpicture}
\end{minipage}
\end{probbox}
\probsol{
\textbf{Understanding the problem.} The spinner has $6$ \emph{equal} sectors, so when it is spun the arrow is equally likely to stop in each sector. We need the probability that the arrow points to a sector marked $P$.
\begin{recallbox}When the outcomes are equally likely,
\[P(\text{event})=\frac{\text{number of favourable outcomes}}{\text{number of all outcomes}}.\]\end{recallbox}
\step{1}{Count all possible outcomes}
The arrow can stop in any one of the $6$ sectors, so there are $6$ possible outcomes.
\step{2}{Count the favourable outcomes}
Reading the letters clockwise, starting from the top-left sector: $P,\ R,\ P,\ Q,\ R,\ Q$. Exactly $2$ sectors are marked $P$ (the top-left one and the right one).
\step{3}{Compute the probability}
$P(\text{letter }P)=\dfrac{2}{6}$ $\left(=\dfrac13\right)$.
\begin{warnbox}There are $3$ different letters ($P$, $Q$, $R$), but the probability is \emph{not} computed as ``$1$ letter out of $3$ letters''. Here $\dfrac13=\dfrac26$ only by coincidence, because each letter appears exactly $2$ times; if the letters took up different numbers of sectors, that way of thinking would be wrong. Always count \emph{sectors}, not \emph{kinds of letters}.\end{warnbox}
\verify{Each of the letters $P$, $Q$, $R$ takes $2$ sectors: $\dfrac26+\dfrac26+\dfrac26=1$ -- the probabilities of all possibilities add up to $1$, as they should.}
The value $\dfrac26$ matches option \textbf{(D)}. (Option (A) $\dfrac36$ comes from miscounting $3$ sectors; (E) $\dfrac16$ is the probability of \emph{one} particular sector.)
\ans{(D) $\dfrac{2}{6}$}
}

% ---------------------------------------------------------------- Problem 2
\begin{probbox}{2}{\relax}
Using the three digits $4$, $7$ and $9$, each digit used \emph{at most once}, how many three-digit whole numbers with different digits can be made?
\choices{$6$}{$3$}{$5$}{$12$}{$9$}
\end{probbox}
\probsol{
\textbf{Understanding the problem.} A three-digit number using exactly the digits $4,7,9$, with no digit repeated. So each number is an \emph{arrangement} of these three digits in the three places: hundreds, tens, ones.
\step{1}{Method 1 -- systematic listing}
Order by the hundreds digit (from smallest to largest):
\begin{itemize}[leftmargin=1.2em,itemsep=0pt,topsep=1pt]
\item Hundreds digit $4$: the other two digits $7,9$ give $479$ and $497$.
\item Hundreds digit $7$: $749$ and $794$.
\item Hundreds digit $9$: $947$ and $974$.
\end{itemize}
There are $2+2+2=6$ numbers.
\step{2}{Method 2 -- multiplication rule}
\begin{recallbox}Multiplication rule: if a task has several successive stages, each with a number of ways (not depending on the choices made before), then the number of ways to do the whole task is the \emph{product} of those numbers.\end{recallbox}
\countboxes{hundreds/3/4 or 7 or 9/tmBlue,tens/2/2 digits left/tmOrange,ones/1/the last digit/tmGreen}{6}
Hundreds: $3$ choices. Tens: only $2$ unused digits remain, $2$ choices. Ones: exactly $1$ digit remains, $1$ choice. Number of numbers: $3\times2\times1=6$.
\begin{warnbox}If repeated digits were allowed there would be $3\times3\times3=27$ numbers (such as $444$, $477$\ldots). The problem forbids repetition, so after each stage the number of choices goes down by $1$.\end{warnbox}
\verify{The six numbers $479,497,749,794,947,974$ are all different and all use the three digits. Both methods give the same result.}
The answer $6$ is option \textbf{(A)}. (Option (B) $3$ is just the number of digits given.)
\ans{(A) $6$}
}

% ---------------------------------------------------------------- Problem 3
\begin{probbox}{3}{\relax}
An envelope contains $2$ sunflower seeds, $3$ mung bean seeds and $4$ pumpkin seeds. Carrie takes one seed from the envelope at random. What is the probability that Carrie takes a sunflower seed?
\choices{$\dfrac{2}{9}$}{$\dfrac{5}{9}$}{$\dfrac{9}{7}$}{$\dfrac{7}{9}$}{$\dfrac{1}{9}$}
\end{probbox}
\probsol{
\textbf{Understanding the problem.} The envelope has three kinds of seeds in different numbers. Carrie takes \emph{one seed} at random: every seed has the same chance of being taken, so each \emph{seed} is one equally likely outcome.
\step{1}{Count all possible outcomes}
Total number of seeds: $2+3+4=9$, i.e.\ $9$ outcomes.
\step{2}{Count the favourable outcomes}
There are $2$ sunflower seeds, so $2$ favourable outcomes.
\step{3}{Compute the probability}
\begin{recallbox}$P=\dfrac{\text{number of favourable outcomes}}{\text{number of all outcomes}}$.\end{recallbox}
$P(\text{sunflower})=\dfrac{2}{9}$.
\begin{warnbox}There are $3$ \emph{kinds} of seeds but the answer is not $\dfrac13$, because the pumpkin seeds are more numerous and so more likely to be taken. Remember too that a probability is never greater than $1$, so option (C) $\dfrac97$ is certainly wrong.\end{warnbox}
\verify{$P(\text{sunflower})+P(\text{mung bean})+P(\text{pumpkin})=\dfrac29+\dfrac39+\dfrac49=\dfrac99=1$.}
The result $\dfrac29$ is option \textbf{(A)}. (Option (D) $\dfrac79$ is the probability of \emph{not} taking a sunflower seed.)
\ans{(A) $\dfrac{2}{9}$}
}

% ---------------------------------------------------------------- Problem 4
\begin{probbox}{4}{\relax}
A cinema has $11$ rows of seats, numbered from $1$ to $11$. Each odd-numbered row has $15$ seats and each even-numbered row has $16$ seats. How many seats are there in the cinema altogether?
\choices{$176$}{$186$}{$165$}{$170$}{$171$}
\end{probbox}
\probsol{
\textbf{Understanding the problem.} The number of seats in a row depends on whether the row number is odd or even. So we need to know how many of the numbers from $1$ to $11$ are odd and how many are even.
\step{1}{Count the odd rows and the even rows}
Odd rows: $1,3,5,7,9,11$ -- $6$ rows.\\
Even rows: $2,4,6,8,10$ -- $5$ rows.\\
Check: $6+5=11$ rows \checkmark.
\begin{recallbox}Number of terms of an evenly spaced sequence $=\dfrac{\text{last}-\text{first}}{\text{step}}+1$. The odd numbers $1,3,\dots,11$ have step $2$: $\dfrac{11-1}{2}+1=6$ numbers. The even numbers $2,4,\dots,10$: $\dfrac{10-2}{2}+1=5$ numbers.\end{recallbox}
\begin{warnbox}Do not split the $11$ rows into ``$5.5$ odd rows and $5.5$ even rows''. The list starts \emph{and} ends with an odd number ($1$ and $11$), so there is exactly $1$ more odd number than even numbers.\end{warnbox}
\step{2}{Seats in the odd rows}
$6\times15=90$ seats.
\step{3}{Seats in the even rows}
$5\times16=80$ seats.
\step{4}{Add up}
Total number of seats: $90+80=170$.
\verify{\emph{Another way -- pairing:} pair row $1$ with row $2$, row $3$ with row $4$, \ldots, row $9$ with row $10$: that gives $5$ pairs, each with $15+16=31$ seats, i.e.\ $5\times31=155$ seats; row $11$ with $15$ seats is left over. Total $155+15=170$ \checkmark.}
The answer $170$ is option \textbf{(D)}. (Mixing up and taking $5$ odd rows and $6$ even rows gives $171$ -- option (E), a very common ``trap''.)
\ans{(D) $170$}
}

% ---------------------------------------------------------------- Problem 5
\begin{probbox}{5}{\relax}
At a class party, each student in turn draws a wrapped present at random from a bag. The presents are only books and calculators; there are $27$ presents in the bag altogether. Meghan draws first. The probability that Meghan draws a book is $\dfrac23$. How many books are in the bag?
\choices{$15$}{$9$}{$21$}{$7$}{$18$}
\end{probbox}
\probsol{
\textbf{Understanding the problem.} A ``backwards'' problem: we know the \emph{probability} and the \emph{total} number of presents, and need the \emph{number of favourable outcomes} (the number of books). Since Meghan draws first, all $27$ presents are still in the bag.
\begin{recallbox}From $P=\dfrac{\text{number of favourable outcomes}}{\text{number of all outcomes}}$ we get: number of favourable outcomes $=P\times$ number of all outcomes.\end{recallbox}
\step{1}{Call the number of books an unknown}
Let the number of books be $s$. The probability of drawing a book is $\dfrac{s}{27}$. By the problem: $\dfrac{s}{27}=\dfrac23$.
\step{2}{Solve the equation}
Multiply both sides of $\dfrac{s}{27}=\dfrac23$ by $27$:
\[s=27\times\frac23=\frac{27\times2}{3}=\frac{54}{3}=18.\]
\emph{Another way -- find the number of calculators first:} let the number of calculators be $c$; then $s+c=27$. The event ``drawing a calculator'' is the complement of ``drawing a book'', so $\dfrac{c}{27}=1-\dfrac23=\dfrac13$, which gives $c=27\times\dfrac13=9$ and $s=27-c=27-9=18$. Both ways give the same result.
\verify{There are $18$ books and $27-18=9$ calculators. Probability of drawing a book: $\dfrac{18}{27}=\dfrac{18\div9}{27\div9}=\dfrac23$ \checkmark.}
\begin{warnbox}The $9$ in option (B) is the number of \emph{calculators}, not books -- read the question carefully before choosing.\end{warnbox}
The answer $18$ is option \textbf{(E)}.
\ans{(E) $18$}
}

% ---------------------------------------------------------------- Problem 6
\begin{probbox}{6}{\relax}
A crate of fruit contains $20$ apples, $10$ oranges and no other fruit. One fruit is taken from the crate at random. What is the probability that it is an orange?
\choices{$\dfrac{1}{10}$}{$\dfrac{1}{20}$}{$\dfrac{1}{30}$}{$\dfrac{1}{3}$}{$\dfrac{2}{3}$}
\end{probbox}
\probsol{
\textbf{Understanding the problem.} Every fruit in the crate has the same chance of being taken. A favourable outcome is taking one of the oranges.
\step{1}{Count all outcomes}
Total number of fruits: $20+10=30$.
\step{2}{Count the favourable outcomes}
There are $10$ oranges.
\step{3}{Compute and simplify}
\begin{recallbox}$P=\dfrac{\text{number of favourable outcomes}}{\text{number of all outcomes}}$; simplify a fraction by dividing the numerator and denominator by their GCD.\end{recallbox}
$P(\text{orange})=\dfrac{10}{30}=\dfrac{10\div10}{30\div10}=\dfrac13$.
\begin{warnbox}Two common mistakes: (1) using the number of apples ($20$) as the denominator instead of the \emph{total number of fruits} ($30$), giving $\dfrac{10}{20}$; (2) writing $\dfrac{1}{10}$ (option A) thinking ``there are $10$ oranges so each is $\dfrac1{10}$''. The denominator is always \emph{all} the outcomes.\end{warnbox}
\verify{$P(\text{apple})=\dfrac{20}{30}=\dfrac23$ and $\dfrac13+\dfrac23=1$ \checkmark. It also makes sense: there are twice as many apples as oranges, so an apple is twice as likely as an orange.}
The result $\dfrac13$ is option \textbf{(D)}. (Option (E) $\dfrac23$ is the probability of taking an apple.)
\ans{(D) $\dfrac13$}
}

% ---------------------------------------------------------------- Problem 7
\begin{probbox}{7}{\relax}
A bag contains $2$ red balls, $5$ yellow balls and $4$ blue balls, all the same size. Without looking, one ball is taken from the bag. The probability of taking a yellow ball is
\choices{$\dfrac{2}{11}$}{$\dfrac{5}{11}$}{$\dfrac{4}{11}$}{$\dfrac{6}{11}$}{$\dfrac{7}{11}$}
\end{probbox}
\probsol{
\textbf{Understanding the problem.} ``Without looking'' and ``the same size'' guarantee that every ball has the same chance of being taken.
\step{1}{Count all outcomes}
Total number of balls: $2+5+4=11$.
\step{2}{Count the favourable outcomes}
There are $5$ yellow balls.
\step{3}{Compute the probability}
\begin{recallbox}$P=\dfrac{\text{number of favourable outcomes}}{\text{number of all outcomes}}$.\end{recallbox}
$P(\text{yellow})=\dfrac{5}{11}$ (this fraction is already in lowest terms because $5$ and $11$ have no common factor other than $1$).
\begin{figC}
\foreach \i in {1,2}{\rvcaBall{\i*0.6}{0}{AnsRed}}
\foreach \i in {3,...,7}{\rvcaBall{\i*0.6+0.3}{0}{StarGold}\draw[StarGold!70!black,line width=0.8pt] (\i*0.6+0.3,0) circle (0.28);}
\foreach \i in {8,...,11}{\rvcaBall{\i*0.6+0.6}{0}{tmBlue}}
\node[font=\footnotesize] at (3.9,-0.55){$5$ of the $11$ balls are yellow};
\end{figC}
\verify{$\dfrac{2}{11}+\dfrac{5}{11}+\dfrac{4}{11}=\dfrac{11}{11}=1$ \checkmark.}
The distractors: (A) $\dfrac{2}{11}$ is the probability of red, (C) $\dfrac{4}{11}$ the probability of blue, (D) $\dfrac{6}{11}$ the probability of \emph{not} taking a yellow ball. The correct answer is \textbf{(B)}.
\ans{(B) $\dfrac{5}{11}$}
}

% ---------------------------------------------------------------- Problem 8
\begin{probbox}{8}{\relax}
You are doing a multiple-choice test. On one question you do not know the answer, so you choose an option at random. If the question has five options (A, B, C, D, E) and only one of them is correct, what is the probability that you guess correctly?
\choices{$\dfrac{1}{5}$}{$\dfrac{5}{5}$}{$\dfrac{4}{5}$}{$\dfrac{2}{5}$}{$\dfrac{3}{5}$}
\end{probbox}
\probsol{
\textbf{Understanding the problem.} ``Choose at random'' means the five options are equally likely to be chosen. The question has exactly \emph{one} correct option.
\step{1}{Count all outcomes}
There are $5$ options to choose from: A, B, C, D, E.
\step{2}{Count the favourable outcomes}
Only $1$ option is correct, so there is $1$ favourable outcome.
\step{3}{Compute the probability}
\begin{recallbox}$P=\dfrac{\text{number of favourable outcomes}}{\text{number of all outcomes}}$.\end{recallbox}
$P(\text{correct})=\dfrac15$.
\begin{tipbox}The probability of guessing \emph{wrong} is $1-\dfrac15=\dfrac45$ -- option (C). So when guessing you are $4$ times more likely to be wrong than right. And (B) $\dfrac55=1$ means ``certainly correct'' -- absurd for a guess.\end{tipbox}
\verify{If you did $100$ such questions and guessed every one, you would get about $100\times\dfrac15=20$ right on average -- matching the feeling ``$1$ right out of every $5$''.}
The answer is \textbf{(A)}.
\ans{(A) $\dfrac15$}
}

% ---------------------------------------------------------------- Problem 9
\begin{probbox}{9}{\relax}
A fair six-sided die has the numbers $1$ to $6$ on its faces. The die is rolled once. What is the probability that it shows a $5$?
\choices{$\dfrac{2}{6}$}{$\dfrac{1}{6}$}{$\dfrac{5}{6}$}{$\dfrac{3}{6}$}{$\dfrac{4}{6}$}
\end{probbox}
\probsol{
\textbf{Understanding the problem.} The die is \emph{fair}, so the six faces are equally likely. The event we care about: ``the face showing is $5$''.
\step{1}{List the possible outcomes}
$\{1,\,2,\,3,\,4,\,5,\,6\}$ -- $6$ equally likely outcomes.
\step{2}{Count the favourable outcomes}
Only one face shows $5$: $1$ favourable outcome.
\step{3}{Compute the probability}
\begin{recallbox}$P=\dfrac{\text{number of favourable outcomes}}{\text{number of all outcomes}}$.\end{recallbox}
$P(5)=\dfrac16$.
\begin{warnbox}Do not confuse ``the number $5$'' with ``the number of favourable outcomes''. Many students write $\dfrac56$ (option C) because they see a $5$ in the question. The $5$ is only the \emph{name} of the face; the number of faces with that name is only $1$. ($\dfrac56$ is the probability of rolling a number \emph{from $1$ to $5$}, i.e.\ anything but $6$.)\end{warnbox}
\verify{Each face has probability $\dfrac16$, and the six faces together give $\dfrac66=1$ \checkmark.}
The answer is \textbf{(B)}.
\ans{(B) $\dfrac16$}
}

% ---------------------------------------------------------------- Problem 10
\begin{probbox}{10}{\relax}
Five friends Andy, Jen, Sally, Mike and Tom sit in a row of five seats. Andy does \emph{not} sit next to Jen, and Sally sits next to Mike. Who \emph{cannot} sit in the middle seat?
\choices{Andy}{Jen}{Sally}{Mike}{Tom}
\end{probbox}
\probsol{
\textbf{Understanding the problem.} The five seats are in a row, numbered $1,2,3,4,5$ from left to right; the middle seat is seat $3$. Two conditions: (i) Andy and Jen do \emph{not} sit next to each other; (ii) Sally and Mike sit next to each other. We must find the friend who can \emph{never} sit in seat $3$.
\begin{recallbox}To show that a friend \emph{can} sit in the middle: just give \emph{one} arrangement satisfying every condition. To show that a friend \emph{cannot}: suppose they sit in the middle and show that every arrangement then breaks a condition (proof by contradiction).\end{recallbox}
\step{1}{Andy, Jen, Sally and Mike can each sit in the middle}
Here is one arrangement for each of them (seats $1\to5$):
\begin{center}
\begin{tabular}{|c|c|c|c|c|c|l|}\hline
In the middle & seat 1 & seat 2 & \textbf{seat 3} & seat 4 & seat 5 & check\\ \hline
Andy & Sally & Mike & \textbf{Andy} & Tom & Jen & Andy (3), Jen (5) not neighbours\\ \hline
Jen & Sally & Mike & \textbf{Jen} & Tom & Andy & Jen (3), Andy (5) not neighbours\\ \hline
Sally & Andy & Mike & \textbf{Sally} & Tom & Jen & Sally (3) next to Mike (2)\\ \hline
Mike & Andy & Sally & \textbf{Mike} & Tom & Jen & Mike (3) next to Sally (2)\\ \hline
\end{tabular}
\end{center}
In all four arrangements Sally and Mike are always neighbours, and Andy and Jen are always apart (Andy in seat $1$ or $3$, Jen in seat $5$, or the other way round). So options (A), (B), (C), (D) are eliminated.
\step{2}{Show that Tom cannot sit in the middle}
Suppose Tom sits in seat $3$. The other four seats split into two pairs of neighbouring seats: the left pair (seats $1$--$2$) and the right pair (seats $4$--$5$). Note that seats $2$ and $4$ are \emph{not} neighbours (Tom is between them).
\begin{figC}
\foreach \i/\n/\c in {1/?/tmBlue,2/?/tmBlue,3/Tom/tmOrange,4/?/tmGreen,5/?/tmGreen}{
  \draw[\c,line width=1pt,fill=\c!12,rounded corners=3pt] (\i*1.5-0.6,0) rectangle (\i*1.5+0.6,0.65);
  \node[font=\bfseries] at (\i*1.5,0.33){\n};\node[font=\footnotesize] at (\i*1.5,-0.28){seat \i};}
\draw[decorate,decoration={brace,amplitude=5pt,mirror},Ink] (0.9,-0.5)--(3.6,-0.5) node[midway,below=5pt,font=\footnotesize]{left pair};
\draw[decorate,decoration={brace,amplitude=5pt,mirror},Ink] (5.4,-0.5)--(8.1,-0.5) node[midway,below=5pt,font=\footnotesize]{right pair};
\end{figC}
\begin{itemize}[leftmargin=1.2em,itemsep=1pt,topsep=1pt]
\item Sally and Mike must be neighbours, and the only neighbouring seats (not counting Tom's) are seats $1$--$2$ or seats $4$--$5$. So Sally and Mike take up one whole pair.
\item Then Andy and Jen are forced into the other pair -- that is, \emph{next to each other}. This contradicts the condition ``Andy does not sit next to Jen''.
\end{itemize}
So the assumption is false: Tom cannot sit in the middle seat.
\begin{warnbox}Do not choose an answer just because ``I tried a few ways and none worked''. For a ``cannot'' question you must argue that \emph{every} way fails, as in Step 2.\end{warnbox}
\verify{You could check by listing all $5\times4\times3\times2\times1=120$ arrangements: no valid one puts Tom in the middle, while each of the other friends has a valid one.}
The answer is \textbf{(E)}.
\ans{(E) Tom}
}

% ---------------------------------------------------------------- Problem 11
\begin{probbox}{11}{\relax}
A box of jelly beans contains $15$ red, $20$ blue and $16$ green jelly beans. Jack takes a green jelly bean and eats it. Then he takes a blue jelly bean and eats it too. If each jelly bean left in the box is equally likely to be chosen, what is the probability that the next jelly bean Jack takes is red?
\choices{$\dfrac{15}{31}$}{$\dfrac{34}{49}$}{$\dfrac{15}{49}$}{$\dfrac{2}{7}$}{$\dfrac{1}{15}$}
\end{probbox}
\probsol{
\textbf{Understanding the problem.} The box has lost $2$ jelly beans before the draw we care about. So we must \emph{recount} the jelly beans in the box just before the third draw.
\begin{recallbox}When objects are taken out and not replaced, update both the \emph{number of favourable outcomes} (numerator) and the \emph{total number of outcomes} (denominator) before computing the probability of the next draw.\end{recallbox}
\step{1}{Jelly beans at the start}
$15+20+16=51$.
\step{2}{Jelly beans after Jack eats $2$}
\begin{center}
\begin{tabular}{|l|c|c|c|c|}\hline
 & red & blue & green & total\\ \hline
At the start & $15$ & $20$ & $16$ & $51$\\ \hline
After eating $1$ green & $15$ & $20$ & $15$ & $50$\\ \hline
After eating $1$ blue & $15$ & $19$ & $15$ & $49$\\ \hline
\end{tabular}
\end{center}
The number of red jelly beans does not change ($15$) because Jack has not eaten a red one; the total goes down to $49$.
\step{3}{Compute the probability}
$P(\text{red})=\dfrac{15}{49}$.
\begin{warnbox}Two common wrong ways: (1) forgetting to subtract the eaten jelly beans, giving $\dfrac{15}{51}$ (not among the options -- a signal to check your work!); (2) comparing the red ones only with the remaining green ones, giving $\dfrac{15}{31}$ (option A). The denominator must be \emph{all} the jelly beans left in the box.\end{warnbox}
\verify{$P(\text{red})+P(\text{not red})=\dfrac{15}{49}+\dfrac{19+15}{49}=\dfrac{15}{49}+\dfrac{34}{49}=1$ \checkmark. (By the way: $\dfrac{34}{49}$ in option (B) is the probability of \emph{not} getting a red one.)}
The answer is \textbf{(C)}.
\ans{(C) $\dfrac{15}{49}$}
}

% ---------------------------------------------------------------- Problem 12
\begin{probbox}{12}{\relax}
Nicholas is counting the sheep in a flock as they cross a road. The sheep start crossing at $2$:$00$ p.m., at a steady rate of $3$ sheep per minute. After counting $42$ sheep, Nicholas falls asleep. He sleeps for exactly an hour and a half and then wakes up, and at that moment the number of sheep that have crossed since $2$:$00$ p.m.\ is exactly half of the whole flock. How many sheep are in the whole flock?
\choices{$630$}{$621$}{$582$}{$624$}{$618$}
\end{probbox}
\probsol{
\textbf{Understanding the problem.} The sheep cross \emph{steadily}, $3$ per minute. The time from $2$:$00$ p.m.\ until Nicholas wakes up has two parts: the time Nicholas spends counting $42$ sheep, and the time Nicholas sleeps ($1$ hour $30$ minutes). The number of sheep that cross in this whole period equals \emph{half the flock}. We let the size of the flock be an unknown and form an equation.
\begin{recallbox}Steady change: amount $=$ rate $\times$ time, so time $=$ amount $\div$ rate.\end{recallbox}
\step{1}{Choose the unknown}
Let the whole flock have $x$ sheep ($x$ a whole number). When Nicholas wakes up, the number of sheep that have crossed since $2$:$00$ p.m.\ is $\dfrac{x}{2}$.
\step{2}{Time Nicholas spends counting $42$ sheep}
$42\div3=14$ minutes. So Nicholas falls asleep at $2$:$14$ p.m.
\step{3}{Total time from $2$:$00$ p.m.\ to waking up}
An hour and a half $=60+30=90$ minutes. Total time: $14+90=104$ minutes, so Nicholas wakes up at $2$:$00$ p.m.\ $+\ 104$ minutes $=3$:$44$ p.m.
\begin{figC}
\draw[->,line width=1pt,Ink] (-0.4,0)--(12.2,0) node[right,font=\footnotesize]{min};
\foreach \x/\m/\t in {0/0/2{:}00,1.54/14/2{:}14,11.44/104/3{:}44}{
  \draw[Ink,line width=1pt] (\x,0.12)--(\x,-0.12);
  \node[below=1pt,font=\footnotesize,align=center] at (\x,-0.12){$\m$\\ $\t$ p.m.};}
\draw[decorate,decoration={brace,amplitude=4pt},tmBlue!80!black,line width=0.8pt] (0,0.22)--(1.54,0.22) node[midway,above=4pt,font=\footnotesize]{counting: $14$ min};
\draw[decorate,decoration={brace,amplitude=5pt},tmOrange!80!black,line width=0.8pt] (1.54,0.22)--(11.44,0.22) node[midway,above=5pt,font=\footnotesize]{asleep: $90$ min};
\end{figC}
\step{4}{Form the equation}
In $104$ minutes, $3\times104=312$ sheep cross. (Or: $42$ sheep while counting $+\ 3\times90=270$ sheep while asleep $=312$ sheep.) This is half the flock, so we have the equation
\[\frac{x}{2}=312.\]
\step{5}{Solve the equation}
Multiply both sides by $2$: $x=312\times2=624$.
\begin{warnbox}Do not forget the $42$ sheep that crossed \emph{before} Nicholas fell asleep. Counting only the $270$ sheep while asleep and doubling gives $540$ -- which does not fit the problem, because it is the number of sheep crossed ``since $2$:$00$ p.m.'' that is half the flock.\end{warnbox}
\verify{Half the flock $=624\div2=312$ sheep, which takes $312\div3=104$ minutes, i.e.\ from $2$:$00$ p.m.\ to $3$:$44$ p.m.; subtracting the $14$ minutes of counting leaves exactly $90$ minutes of sleep \checkmark.}
The answer is \textbf{(D)}.
\ans{(D) $624$}
}

% ---------------------------------------------------------------- Problem 13
\needspace{10\baselineskip}
\begin{probbox}{13}{\relax}
A positive whole number is called a \emph{palindrome} if its digits read the same from left to right as from right to left. For example, $474$ and $222$ are palindromes. How many palindromes are there between $100$ and $1000$?
\choices{$10$}{$90$}{$100$}{$900$}{$1000$}
\end{probbox}
\probsol{
\textbf{Understanding the problem.} The numbers between $100$ and $1000$ are the \emph{three-digit} numbers ($101$ to $999$; $100$ and $1000$ are not palindromes, so including them or not changes nothing). A three-digit number $\overline{abc}$ is a palindrome when reading it backwards, $\overline{cba}$, gives the same number, that is $c=a$. So a palindrome has the form $\overline{aba}$.
\step{1}{Choose the hundreds digit $a$}
$a$ is the first digit, so $a\neq0$: $a\in\{1,2,\dots,9\}$ -- $9$ ways.
\step{2}{Choose the tens digit $b$}
$b$ can be any digit from $0$ to $9$ -- $10$ ways.
\step{3}{The ones digit}
The ones digit must equal $a$ -- only $1$ way (it cannot be chosen freely).
\step{4}{Apply the multiplication rule}
\begin{recallbox}Multiplication rule: number of ways $=$ product of the numbers of choices at each stage.\end{recallbox}
\countboxes{hundreds a/9/from 1 to 9/tmBlue,tens b/10/from 0 to 9/tmOrange,ones/1/must equal a/tmGreen}{90}
There are $9\times10\times1=90$ palindromes.
\verify{List by $a$: for $a=1$ we have $101,111,121,\dots,191$ -- exactly $10$ numbers. Each value of $a$ from $1$ to $9$ gives $10$ such numbers, so the total is $9\times10=90$ \checkmark.}
\begin{warnbox}Option (D) $900$ is the number of \emph{all} three-digit numbers (from $100$ to $999$), not the palindromes. Option (C) $100$ comes from forgetting the condition $a\neq 0$ (counting ``$010$'', ``$020$''\ldots, which are not three-digit numbers).\end{warnbox}
The answer is \textbf{(B)}.
\ans{(B) $90$}
}

% ---------------------------------------------------------------- Problem 14
\begin{probbox}{14}{\relax}
\begin{minipage}[t]{0.5\linewidth}\vspace{0pt}
The figure shown is made only of octagons and small squares. The ratio of the number of octagons to the number of squares is
\choicesII{$1:1$}{$2:1$}{$25:12$}{$5:0$}{$5:3$}
\end{minipage}\hfill
\begin{minipage}[t]{0.46\linewidth}\vspace{0pt}
\centering\begin{tikzpicture}[scale=0.62]
\foreach \i in {0,...,4}{\foreach \j in {0,...,3}{\cpOct{\i}{\j}}}
\foreach \i in {0.5,1.5,2.5,3.5}{\foreach \j in {0.5,1.5,2.5}{\cpDia{\i}{\j}}}
\end{tikzpicture}
\end{minipage}
\end{probbox}
\probsol{
\textbf{Understanding the problem.} We count the octagons and the squares separately, then form the ratio (octagons : squares) and simplify.
\step{1}{Count the octagons}
The octagons are arranged in $4$ rows of $5$. Number of octagons: $4\times5=20$.
\step{2}{Count the squares}
Each small square sits where \emph{four} octagons meet (between two neighbouring rows and between two neighbouring columns).
\begin{itemize}[leftmargin=1.2em,itemsep=0pt,topsep=1pt]
\item Between $5$ columns of octagons there are $5-1=4$ vertical gaps.
\item Between $4$ rows of octagons there are $4-1=3$ horizontal gaps.
\end{itemize}
Number of squares: $3\times4=12$ (exactly what you see when counting directly in the figure: $3$ rows of $4$).
\begin{warnbox}There are \emph{no} squares on the outer edge of the figure (the notch between two octagons on the edge is not closed). Do not count extra ``half squares'' around the border.\end{warnbox}
\step{3}{Form and simplify the ratio}
\begin{recallbox}Simplify a ratio $a:b$ by dividing both numbers by their GCD.\end{recallbox}
Ratio: $20:12$. GCD$(20,12)=4$, so $20:12=(20\div4):(12\div4)=5:3$.
\verify{$5:3$ means $3$ squares for every $5$ octagons; $20=5\times4$ and $12=3\times4$ \checkmark.}
The distractors: (D) $5:0$ forgets to count the squares; (C) $25:12$ miscounts the octagons as $5\times5$. The answer is \textbf{(E)}.
\ans{(E) $5:3$}
}

% ---------------------------------------------------------------- Problem 15
\begin{probbox}{15}{\relax}
\begin{minipage}[t]{0.52\linewidth}\vspace{0pt}
The figure shown consists of three identical rectangles overlapping each other. How many rectangles of all sizes are there in the figure?
\choicesII{$11$}{$15$}{$7$}{$13$}{$9$}
\end{minipage}\hfill
\begin{minipage}[t]{0.44\linewidth}\vspace{0pt}
\centering\begin{tikzpicture}[x=0.95cm,y=0.42cm]
\draw[Ink,line width=1pt] (0,2) rectangle (2,4);
\draw[Ink,line width=1pt] (1,1) rectangle (3,3);
\draw[Ink,line width=1pt] (2,0) rectangle (4,2);
\end{tikzpicture}
\end{minipage}
\end{probbox}
\probsol{
\textbf{Understanding the problem.} Call the three identical rectangles \textbf{I} (top left), \textbf{II} (middle) and \textbf{III} (bottom right). Where their sides cross, more small rectangles appear. We must count \emph{all} rectangles whose four sides are all drawn in the figure.
\step{1}{Look carefully at the figure}
Label the points as below. Notice: the bottom side of rectangle I and the top side of rectangle III lie on \emph{the same straight line} and join at the point $O$ -- the centre of rectangle II. Likewise the right side of I and the left side of III join at $O$. So inside rectangle II there is a ``plus sign'' that divides II into $4$ equal small cells.
\begin{center}\begin{tikzpicture}[x=1.5cm,y=0.75cm,font=\sffamily\small]
\fill[tmOrange!25] (1,2) rectangle (2,3);\fill[tmGreen!25] (2,2) rectangle (3,3);
\fill[tmBlue!20] (1,1) rectangle (2,2);\fill[tmPurple!20] (2,1) rectangle (3,2);
\draw[Ink,line width=1pt] (0,2) rectangle (2,4);
\draw[Ink,line width=1pt] (1,1) rectangle (3,3);
\draw[Ink,line width=1pt] (2,0) rectangle (4,2);
\node at (1.5,2.5){$1$};\node at (2.5,2.5){$2$};\node at (1.5,1.5){$3$};\node at (2.5,1.5){$4$};
\fill[AnsRed] (2,2) circle (2pt);\node[above right,font=\footnotesize\bfseries,text=AnsRed] at (2,2){$O$};
\node[font=\bfseries,text=Slate] at (0.5,3.4){I};\node[font=\bfseries,text=Slate] at (3.5,0.6){III};
\node[font=\bfseries,text=Slate,anchor=west] at (3.05,2.75){II};
\end{tikzpicture}\end{center}
\step{2}{Count type by type}
\begin{itemize}[leftmargin=1.2em,itemsep=1pt,topsep=1pt]
\item \textbf{The three original large rectangles:} I, II, III -- $3$ rectangles.
\item \textbf{The four small cells} inside rectangle II (cells $1$, $2$, $3$, $4$) -- $4$ rectangles. Each cell has all four sides: for example cell $1$ has its top and left sides on rectangle II, and its bottom and right sides on rectangle I.
\item \textbf{Halves of rectangle II} (two neighbouring cells joined): cells $1+2$ (top half), $3+4$ (bottom half), $1+3$ (left half), $2+4$ (right half) -- $4$ rectangles.
\end{itemize}
\step{3}{Check that there are no others}
\begin{itemize}[leftmargin=1.2em,itemsep=1pt,topsep=1pt]
\item The part of rectangle I \emph{outside} rectangle II is L-shaped, and no segment passes through it. A rectangle containing part of this L-shaped region must use the left side and the top side of rectangle I as its sides; these two sides are exactly as long as the sides of rectangle I, so that rectangle is rectangle I itself. Likewise, a rectangle containing part of III outside II can only be rectangle III.
\item Every other rectangle must lie entirely inside rectangle II. Rectangle II with its inner ``plus sign'' is a $2\times2$ grid of cells, and in Step 2 we have listed them all: $4$ cells, $4$ halves and rectangle II itself.
\end{itemize}
\begin{warnbox}Do not count the large rectangle surrounding all three shapes (from the top-left corner of I to the bottom-right corner of III): \emph{none} of its sides is fully drawn. And do not miss the small cells created by the ``plus sign'' in the centre.\end{warnbox}
\step{4}{Add up}
$3+4+4=11$ rectangles.
\verify{Another count -- inside rectangle II alone (a complete $2\times2$ grid): by Example 7 there are $9$ rectangles ($4$ cells, $4$ halves and II itself). Adding rectangles I and III: $9+2=11$ \checkmark.}
The answer is \textbf{(A)}.
\ans{(A) $11$}
}
